\documentclass[11pt]{article}
\usepackage[margin=2.5cm]{geometry}
\usepackage{amsmath,amssymb}
\usepackage{booktabs}
\usepackage{hyperref}

\title{Exact Coordinates for Heilbronn Configurations\\[4pt]
       \large $n = 5, \ldots, 10$ on the Unit Square}
\author{\parbox{\dimexpr\textwidth-2\tabcolsep}{\centering\normalsize Companion to: N.\ Sudermann-Merx, \emph{From Global Optimization to Exact Coordinates:
Heilbronn's Triangle Problem on the Unit Square Using Mixed-Integer Nonlinear Optimization}, 2026.}}
\date{}

\begin{document}
\maketitle

\noindent
This document lists exact coordinates of configurations for
Heilbronn's triangle problem on the unit square~$[0,1]^2$
for $n = 5, \ldots, 10$.
They were obtained by solving a mixed-integer nonlinear program (MINLP)
with Gurobi to global optimality within solver tolerances and then deriving
exact symbolic coordinates with SymPy. All triangle areas are verified in exact
arithmetic. The configurations are proven optimal for $n\le 8$ (see the paper
for references). For $n=9$ and $n=10$, the value is a rigorous lower bound that matches
the solver's numerical upper bound, but no rigorous proof of optimality is known.
For $n=10$, the upper bound also rests on the numerical upper bound for $\Delta_9$.
The derivation is fully automated by the script
\texttt{exact\_coordinates.py} included in this repository.

Points are labeled $1, \ldots, n$.
We denote the minimum triangle area of the listed configuration by~$\Delta(P)$, which equals $\Delta_n$ for $n\le 8$.

%% ===================================================================
\section*{$n = 5$: \quad $\Delta_5 = \dfrac{\sqrt{3}}{9} \approx 0.19245$}

Four critical triangles. The polynomial system of the critical triangles has a unique solution in the unit square.

\begin{center}
\begin{tabular}{c c c}
\toprule
Pt & $x$ & $y$ \\
\midrule
$1$ & $0$ & $\frac{1}{3}$ \\[4pt]
$2$ & $\frac{\sqrt{3}}{3}$ & $0$ \\[4pt]
$3$ & $1$ & $1 - \frac{\sqrt{3}}{3}$ \\[4pt]
$4$ & $\frac{2}{3}$ & $1$ \\[4pt]
$5$ & $0$ & $1$ \\
\bottomrule
\end{tabular}
\end{center}

%% ===================================================================
\section*{$n = 6$: \quad $\Delta_6 = \dfrac{1}{8} = 0.125$}

Six critical triangles. One-parameter family of optimal solutions
parameterized by $c \in [0, \tfrac{1}{4}]$.

\begin{center}
\begin{tabular}{c c c}
\toprule
Pt & $x$ & $y$ \\
\midrule
$1$ & $0$ & $c$ \\[4pt]
$2$ & $\frac{1}{2}$ & $0$ \\[4pt]
$3$ & $1$ & $\frac{1}{2} - c$ \\[4pt]
$4$ & $\frac{1}{2}$ & $1$ \\[4pt]
$5$ & $0$ & $c + \frac{1}{2}$ \\[4pt]
$6$ & $1$ & $1-c$ \\
\bottomrule
\end{tabular}
\end{center}

%% ===================================================================
\section*{$n = 7$: \quad $\Delta_7 = -\dfrac{a}{2a-4} \approx 0.08386$}

Eight critical triangles. Coordinates are rational functions of
$a \approx 0.2873$, the unique root in $(0,1)$ of the irreducible cubic
\[
  a^3 + 5\,a^2 - 5\,a + 1 = 0.
\]

\begin{center}
\begin{tabular}{c c c}
\toprule
Pt & $x$ & $y$ \\
\midrule
$1$ & $0$ & $0$ \\[4pt]
$2$ & $\dfrac{1}{a^2-3a+2}$ & $0$ \\[8pt]
$3$ & $1$ & $a$ \\[4pt]
$4$ & $\dfrac{2a-1}{a(a-2)}$ & $1$ \\[8pt]
$5$ & $0$ & $1$ \\[4pt]
$6$ & $\dfrac{a-1}{a-2}$ & $a$ \\[8pt]
$7$ & $\dfrac{2(-a^2+3a-1)}{a^3-4a^2+5a-2}$ & $-\dfrac{2a}{a-1}$ \\[8pt]
\bottomrule
\end{tabular}
\end{center}

%% ===================================================================
\section*{$n = 8$: \quad $\Delta_8 = -\dfrac{1}{36}+\dfrac{\sqrt{13}}{36} \approx 0.07238$}

Twelve critical triangles. Point-symmetric configuration with
the polynomial system of the critical triangles has a unique solution in the unit square, involving $\sqrt{13}$ (this does not imply global uniqueness).

\begin{center}
\begin{tabular}{c c c}
\toprule
Pt & $x$ & $y$ \\
\midrule
$1$ & $0$ & $0$ \\[4pt]
$2$ & $\frac{1}{6}+\frac{\sqrt{13}}{6}$ & $0$ \\[4pt]
$3$ & $1$ & $\frac{7}{18}-\frac{\sqrt{13}}{18}$ \\[4pt]
$4$ & $1$ & $1$ \\[4pt]
$5$ & $0$ & $\frac{11}{18}+\frac{\sqrt{13}}{18}$ \\[4pt]
$6$ & $\frac{5}{6}-\frac{\sqrt{13}}{6}$ & $1$ \\[4pt]
$7$ & $\frac{5}{6}-\frac{\sqrt{13}}{6}$ & $\frac{7}{9}-\frac{\sqrt{13}}{9}$ \\[4pt]
$8$ & $\frac{1}{6}+\frac{\sqrt{13}}{6}$ & $\frac{2}{9}+\frac{\sqrt{13}}{9}$ \\
\bottomrule
\end{tabular}
\end{center}

%% ===================================================================
\section*{$n = 9$: \quad $\Delta(P) = -\dfrac{11}{64}+\dfrac{9\sqrt{65}}{320} \approx 0.05488$}

Eleven critical triangles. Configuration symmetric under the reflection
in a diagonal. The polynomial system of the critical triangles has a unique solution in the unit square, involving $\sqrt{65}$ (this does not imply global uniqueness).

\begin{center}
\begin{tabular}{c c c}
\toprule
Pt & $x$ & $y$ \\
\midrule
$1$ & $0$ & $1-\frac{\sqrt{65}}{10}$ \\[4pt]
$2$ & $\frac{3}{8}-\frac{\sqrt{65}}{40}$ & $0$ \\[4pt]
$3$ & $1$ & $\frac{9}{16}-\frac{3\sqrt{65}}{80}$ \\[4pt]
$4$ & $\frac{3}{8}-\frac{\sqrt{65}}{40}$ & $1$ \\[4pt]
$5$ & $0$ & $\frac{5}{8}+\frac{\sqrt{65}}{40}$ \\[4pt]
$6$ & $\frac{1}{4}+\frac{\sqrt{65}}{20}$ & $\frac{3}{4}-\frac{\sqrt{65}}{20}$ \\[4pt]
$7$ & $\frac{7}{16}+\frac{3\sqrt{65}}{80}$ & $0$ \\[4pt]
$8$ & $\frac{\sqrt{65}}{10}$ & $1$ \\[4pt]
$9$ & $1$ & $\frac{5}{8}+\frac{\sqrt{65}}{40}$ \\
\bottomrule
\end{tabular}
\end{center}

%% ===================================================================
\section*{$n = 10$: \quad $\Delta(P) = \dfrac{5}{8}g^2-\dfrac{1}{2}g^3 \approx 0.04654$}

Sixteen critical triangles. Configuration symmetric under the reflections in
both diagonals, found by Comellas and Yebra (2002). Here $g\approx 0.3156$ is the real root of
$12g^3-27g^2+20g-4$, explicitly
$g=\frac34-\frac{(63+8\sqrt{62})^{1/3}}{12}-\frac{1}{12(63+8\sqrt{62})^{1/3}}$.
The equal-area conditions alone leave a one-parameter family. The configuration
is the point of this family with the largest common area, obtained with an
additional stationarity condition.

\begin{center}
\begin{tabular}{c c c}
\toprule
Pt & $x$ & $y$ \\
\midrule
$1$ & $0$ & $1-3g+2g^2$ \\[4pt]
$2$ & $1-3g+2g^2$ & $0$ \\[4pt]
$3$ & $1$ & $3g-2g^2$ \\[4pt]
$4$ & $3g-2g^2$ & $1$ \\[4pt]
$5$ & $0$ & $1-\frac{g}{2}$ \\[4pt]
$6$ & $\frac{g}{2}$ & $1$ \\[4pt]
$7$ & $g$ & $g$ \\[4pt]
$8$ & $1-g$ & $1-g$ \\[4pt]
$9$ & $1-\frac{g}{2}$ & $0$ \\[4pt]
$10$ & $1$ & $\frac{g}{2}$ \\
\bottomrule
\end{tabular}
\end{center}

\bigskip\noindent
\textbf{Source:}
\url{https://github.com/spiralulam/heilbronn}

\end{document}
